\documentclass[11pt]{article}
\usepackage[margin=1in]{geometry}
\usepackage{amsmath,amssymb}
\usepackage[hidelinks]{hyperref}
\title{Literature resolution: cyclic polynomials on the bidisk}
\author{Run \texttt{fa\_banach\_001}, lane 13}
\date{9 August 2026}
\begin{document}
\maketitle

\paragraph{Source question.}
Problem 5.2 of arXiv:1310.4094 asks for a characterization, for every
$\alpha\leq 1$, of the cyclic polynomials in the product-weight
Dirichlet-type space $\mathfrak D_\alpha$ on the bidisk.

\paragraph{Later answer.}
ArXiv:1408.3616 explicitly identifies Problem 5.2 and states that it solves
the problem. Its Main Theorem says that if $f\in\mathbb C[z_1,z_2]$ is
irreducible and has no zero in $\mathbb D^2$, then:
\begin{enumerate}
\item if $\alpha\leq \tfrac12$, then $f$ is cyclic;
\item if $\tfrac12<\alpha\leq1$, then $f$ is cyclic exactly when
$\mathcal Z(f)\cap\mathbb T^2$ is empty or finite, or $f$ is a constant
multiple of $\zeta-z_1$ or $\zeta-z_2$ for some $\zeta\in\mathbb T$;
\item if $\alpha>1$, then $f$ is cyclic exactly when
$\mathcal Z(f)\cap\mathbb T^2$ is empty.
\end{enumerate}
Since polynomials are multipliers, a general polynomial is cyclic exactly
when each irreducible factor is cyclic. A zero in the open bidisk always
precludes cyclicity.

\paragraph{Identification.}
The two papers use the identical coefficient norm
\[
 \lVert f\rVert_\alpha^2=\sum_{k,l\geq0}(k+1)^\alpha(l+1)^\alpha
 |a_{k,l}|^2,
\]
so the later Main Theorem is an exact answer, not an analogous result for a
different scale of spaces. The Brown--Shields question for arbitrary outer
functions remains outside this packet.

\paragraph{References.}
C. Beneteau, A. A. Condori, C. Liaw, D. Seco, and A. A. Sola,
\emph{Cyclicity in Dirichlet-type spaces and extremal polynomials II:
functions on the bidisk}, arXiv:1310.4094.

C. Beneteau, G. Knese, L. Kosinski, C. Liaw, D. Seco, and A. Sola,
\emph{Cyclic polynomials in two variables}, arXiv:1408.3616.
\end{document}
